\documentclass[11pt,a4paper]{article}

\usepackage[T1]{fontenc}
\usepackage{lmodern}
\usepackage[margin=22mm]{geometry}
\usepackage{amsmath,amssymb,amsthm}
\usepackage{microtype}
\usepackage[hidelinks]{hyperref}

\hypersetup{
  pdftitle={Prescribed even cycle lengths in cubic graphs},
  pdfsubject={A girth-dependent sufficient condition, with a full proof and a power-of-two corollary}
}
\newtheorem{theorem}{Theorem}
\newtheorem{corollary}[theorem]{Corollary}
\DeclareMathOperator{\tr}{tr}
\DeclareMathOperator{\head}{head}
\DeclareMathOperator{\tail}{tail}
\setlength{\parindent}{1.2em}
\setlength{\parskip}{0pt}
\allowdisplaybreaks[1]
\raggedbottom

% Based on the nonbacktracking-walk counting argument in the supplied
% draft cubic-girth-power-of-two.pdf and its previously expanded proof.
% The target length is any admissible even integer, not necessarily
% a power of two. No connectedness or expansion assumption is imposed.
% The spectral-refinement and forest addenda are not used.
% This file is standalone and can be compiled with pdflatex (two passes).

\begin{document}

\begin{center}
  {\Large\bfseries Prescribed even cycle lengths in cubic graphs}
\end{center}
\vspace{0.2em}

\begin{theorem}[Prescribed even lengths]\label{thm:main}
Let $G$ be a finite simple cubic graph on $n\ge64$ vertices, with girth
$g$, and put
\[
  h:=\left\lfloor\frac{g-1}{2}\right\rfloor.
\]
For every even integer $\ell$ satisfying
\begin{equation}\label{eq:hypotheses}
  \ell\ge 2\log_2 n+4
  \qquad\text{and}\qquad
  2^{2h}>\frac92\,n\ell^2,
\end{equation}
$G$ contains a simple cycle of length exactly $\ell$.
Equivalently, if $C(G)$ is the set of lengths of simple cycles in $G$,
then
\begin{equation}\label{eq:interval}
  \left[\,2\log_2 n+4,\;
    \frac{\sqrt2}{3}\,\frac{2^h}{\sqrt n}\,\right)
  \cap 2\mathbb Z
  \;\subseteq\; C(G).
\end{equation}
The graph need not be connected. All logarithms displayed with
subscript $2$ are to base $2$.
\end{theorem}

\begin{proof}
Fix an even integer $\ell$ satisfying \eqref{eq:hypotheses}.
Write $G=(V,E)$ and $m:=|E|=3n/2$. A cycle is simple, and its length
is its number of edges. The definition of $h$ gives
\begin{equation}\label{eq:h}
  1\le h<g,
  \qquad g-2\le 2h<g.
\end{equation}

\medskip
\noindent\textbf{1. Counting strictly nonbacktracking closed walks.}
Let $\vec E$ be the set of $2m=3n$ directed edges obtained by giving
each edge of $G$ its two orientations. For $e\in\vec E$, let $\bar e$
be its reverse. Define the nonbacktracking matrix $B$, with rows and
columns indexed by $\vec E$, by
\[
  B_{e,f}:=
  \begin{cases}
    1,&\text{if }\head(e)=\tail(f)\text{ and }f\ne\bar e,\\
    0,&\text{otherwise}.
  \end{cases}
\]
Set $T_\ell:=\tr(B^\ell)$. Expanding the trace gives
\begin{equation}\label{eq:trace-count}
  T_\ell=
  \sum_{e_0,\ldots,e_{\ell-1}\in\vec E}
    \prod_{j=0}^{\ell-1}B_{e_j,e_{j+1}},
  \qquad e_\ell:=e_0.
\end{equation}
Thus $T_\ell$ counts ordered directed-edge sequences describing
closed walks of length $\ell$ with no immediate reversal, including
between the last and first edges. These will be called
\emph{rooted oriented strictly nonbacktracking closed walks}.
The first directed edge fixes the root and the direction of traversal;
sequences are not identified under rotation or reversal. In vertex
notation, such a sequence is $(x_0,\ldots,x_\ell=x_0)$ with
$x_{j-1}\ne x_{j+1}$ for every $j$ read modulo $\ell$.

Let $A$ be the adjacency matrix of $G$. We will use the identity
\begin{equation}\label{eq:bass}
  \det(zI_{2m}-B)
  =(z^2-1)^{m-n}\det\bigl(z^2I_n-zA+2I_n\bigr).
\end{equation}
This is the cubic case of the Ihara--Bass determinant identity;
we derive it here.\footnote{For the general identity and an
incidence-matrix proof, see M.~Kempton, \emph{Non-backtracking random walks
and a weighted Ihara's theorem}, arXiv:1603.05553 (2016), \S\,2.4.}
Define $n\times2m$ matrices $S,H$ and a $2m\times2m$ matrix $J$ by
\[
  S_{v,e}=\mathbf1_{\{v=\tail(e)\}},
  \qquad
  H_{v,e}=\mathbf1_{\{v=\head(e)\}},
  \qquad
  J_{e,f}=\mathbf1_{\{f=\bar e\}},
\]
where $\mathbf1_{\{\cdot\}}$ is the indicator of the stated condition.
Directly from these definitions and the fact that every degree is $3$,
\[
  B=H^{\mathsf T}S-J,
  \qquad SH^{\mathsf T}=A,
  \qquad SJH^{\mathsf T}=3I_n,
  \qquad J^2=I_{2m}.
\]
Since $J$ exchanges the two orientations of each edge, for $z\ne\pm1$
we have
\[
  \det(zI_{2m}+J)=(z^2-1)^m,
  \qquad
  (zI_{2m}+J)^{-1}=\frac{zI_{2m}-J}{z^2-1}.
\]
For an invertible square matrix $M$ and compatible rectangular
matrices $U,V$, the determinant formula
\[
  \det(M-UV)=\det(M)\det(I-VM^{-1}U)
\]
follows by eliminating either off-diagonal block in
$\begin{pmatrix}M&U\\V&I\end{pmatrix}$.
Applying it with $M=zI_{2m}+J$, $U=H^{\mathsf T}$ and $V=S$ yields
\begin{align*}
  \det(zI_{2m}-B)
  &=\det\bigl(zI_{2m}+J-H^{\mathsf T}S\bigr)\\
  &=(z^2-1)^m
    \det\!\left(I_n-S(zI_{2m}+J)^{-1}H^{\mathsf T}\right)\\
  &=(z^2-1)^m
    \det\!\left(I_n-\frac{zA-3I_n}{z^2-1}\right)\\
  &=(z^2-1)^{m-n}\det\bigl(z^2I_n-zA+2I_n\bigr).
\end{align*}
Since $m-n=n/2\ge0$, both sides are polynomials in $z$; equality for
$z\ne\pm1$ therefore proves \eqref{eq:bass} for every $z$.
No connectedness assumption has been used.

Let $\lambda_1,\ldots,\lambda_n$ be the real eigenvalues of the
symmetric matrix $A$, counted with multiplicity. Because
$A\mathbf1=3\mathbf1$, take $\lambda_1=3$; this remains valid for a
disconnected graph. Formula \eqref{eq:bass} shows that the eigenvalues
of $B$, with algebraic multiplicity, consist of the roots
$\mu_{i,+},\mu_{i,-}$ of
\[
  \mu^2-\lambda_i\mu+2=0
  \qquad(1\le i\le n),
\]
together with $m-n$ additional copies of each of $1$ and $-1$.
For any square matrix, triangularization over $\mathbb C$ shows that
the trace of its $\ell$th power is the sum of the $\ell$th powers of
its eigenvalues, with algebraic multiplicity. In particular, this
assertion does not require $B$ to be diagonalizable.

The roots associated with $\lambda_1=3$ are $2$ and $1$.
If $|\lambda_i|\ge2\sqrt2$, its two roots are real, and their
$\ell$th powers are nonnegative because $\ell$ is even. If
$|\lambda_i|<2\sqrt2$, the roots are complex conjugates with modulus
$\sqrt2$, so
\[
  \mu_{i,+}^{\,\ell}+\mu_{i,-}^{\,\ell}
  =2\operatorname{Re}\!\left(\mu_{i,+}^{\,\ell}\right)
  \ge-2\,2^{\ell/2}.
\]
The additional eigenvalues $1$ and $-1$ also contribute
nonnegatively. Keeping the contribution $2^\ell$ from the first pair
and bounding each remaining pair from below gives
\begin{equation}\label{eq:trace-spectral}
  T_\ell\ge 2^\ell-2(n-1)2^{\ell/2}
  >2^\ell-2n\,2^{\ell/2}.
\end{equation}
Since $\ell\ge2\log_2 n+4$, we have $2^{\ell/2}\ge4n$, whence
\[
  2n\,2^{\ell/2}\le2^{\ell-1}.
\]
Consequently,
\begin{equation}\label{eq:trace-lower}
  T_\ell>2^{\ell-1}.
\end{equation}

\medskip
\noindent\textbf{2. Bounding internally nonbacktracking returns.}
For a vertex $v$ and an integer $r\ge1$, let $f_r(v)$ count the
ordered vertex sequences
\[
  (y_0,y_1,\ldots,y_r),
  \qquad y_0=y_r=v,
\]
whose consecutive vertices are adjacent and which satisfy
\[
  y_{j-1}\ne y_{j+1}\qquad(1\le j\le r-1).
\]
These are \emph{internally nonbacktracking returns}; no condition is
imposed between the last edge and the first edge. We claim that
\begin{equation}\label{eq:return-bound}
  f_r(v)=0\quad(1\le r<g),
  \qquad
  f_r(v)\le\frac32\,2^{r-h}\quad(r\ge g).
\end{equation}

Indeed, every internally nonbacktracking walk that repeats a vertex
contains a simple cycle. To see this, choose the smallest index $j$
at which a vertex repeats, and let $i<j$ be its earlier occurrence.
All vertices preceding index $j$ are distinct. Hence the segment
from index $i$ to index $j$ has distinct vertices apart from its
identical endpoints. Its length is neither $1$, which would be a
loop, nor $2$, which would be an immediate reversal. It is therefore
a simple cycle. It follows that every internally nonbacktracking
walk of length less than $g$ is a simple path, and in particular no
positive-length return of length less than $g$ exists.
This proves the first assertion in \eqref{eq:return-bound}.

For $r\ge g$, split each return counted by $f_r(v)$ into a prefix of
length $r-h$ and a suffix of length $h$. There are exactly
\[
  3\cdot2^{r-h-1}
\]
internally nonbacktracking walks of length $r-h$ starting at $v$:
there are three choices at the first step and two at each later
step. Here $r-h\ge1$ by \eqref{eq:h}.
Fix such a prefix and let its endpoint be $w$.
Every admissible suffix is a simple path from $w$ to $v$, since
$h<g$. If $w=v$, no such positive-length simple path exists.
If $w\ne v$, there is at most one such path of length $h$:
the union of two distinct simple paths between the same endpoints
contains a cycle and has at most $2h<g$ edges, contradicting the
girth. Any nonbacktracking requirement at the prefix--suffix join
can only remove possibilities. Each prefix thus has at most one
admissible suffix, proving
\[
  f_r(v)\le3\cdot2^{r-h-1}
  =\frac32\,2^{r-h}.
\]

\medskip
\noindent\textbf{3. Bounding the nonsimple strictly nonbacktracking walks.}
Let $D_\ell$ count the walks in $T_\ell$ for which a vertex repeats
among $x_0,\ldots,x_{\ell-1}$. Let $c_\ell$ be the number of simple
cycles of length $\ell$, identified under cyclic rotation and
reversal. Every such cycle has exactly $\ell$ distinct choices of
starting vertex and two choices of direction. Thus
\begin{equation}\label{eq:cycle-count}
  T_\ell=2\ell\,c_\ell+D_\ell.
\end{equation}

Consider a walk counted by $D_\ell$. Choose indices
$0\le i<j\le\ell-1$ with $x_i=x_j$, and extend its vertex sequence
periodically by $x_{k+\ell}=x_k$ for all integers $k$.
Set $v=x_i$ and $r=j-i$.
The two segments
\[
  R_1=(x_i,x_{i+1},\ldots,x_j),
  \qquad
  R_2=(x_j,x_{j+1},\ldots,x_{i+\ell})
\]
are returns to $v$ of positive lengths $r$ and $\ell-r$.
Both are internally nonbacktracking: every internal consecutive
edge pair is a consecutive pair in the original cyclic walk.
In particular, the original last--first pair remains nonbacktracking
when it lies inside $R_2$. By \eqref{eq:return-bound},
\begin{equation}\label{eq:split-range}
  g\le r\le\ell-g.
\end{equation}
The artificial closing pair of either return may reverse, which is
permitted by the definition of $f_r(v)$.

For each fixed $v$ and $r$, there are at most
$f_r(v)f_{\ell-r}(v)$ ordered pairs $(R_1,R_2)$.
A pair determines a concatenated closed vertex sequence.
Placing its starting point at one of the $\ell$ edge positions
produces at most $\ell$ rooted oriented sequences, and every walk
counted by $D_\ell$ is produced by some such pair and position.
The orientation is already specified by the ordered returns, so
there is no additional factor of two. Some concatenations have
inadmissible joins, and some walks have multiple representations;
retaining all of them can only enlarge this upper bound. For a
periodic walk, different starting positions can give the same
sequence, which is why the number of resulting sequences is
\emph{at most} $\ell$.
Consequently,
\begin{align}
  D_\ell
  &\le\ell\sum_{r=g}^{\ell-g}\sum_{v\in V}
       f_r(v)f_{\ell-r}(v)\notag\\
  &\le\ell\cdot\ell\cdot n\cdot\frac94\,2^{\ell-2h}
   =\frac94\,n\ell^2\,2^{\ell-2h}.
  \label{eq:nonsimple-bound}
\end{align}
If $\ell<2g$, \eqref{eq:split-range} is impossible, so $D_\ell=0$
and the sum is empty. Otherwise there are at most $\ell$ possible
integer values of $r$, and \eqref{eq:return-bound} bounds each
product in the sum by
\[
  \left(\frac32\,2^{r-h}\right)
  \left(\frac32\,2^{\ell-r-h}\right)
  =\frac94\,2^{\ell-2h}.
\]
Thus \eqref{eq:nonsimple-bound} holds in all cases.

\medskip
\noindent\textbf{4. Extracting a simple cycle.}
Combining \eqref{eq:trace-lower}, \eqref{eq:cycle-count} and
\eqref{eq:nonsimple-bound}, we obtain
\begin{align*}
  2\ell\,c_\ell
  &=T_\ell-D_\ell\\
  &>2^{\ell-1}-\frac94\,n\ell^2\,2^{\ell-2h}\\
  &=2^{\ell-1}\left(1-\frac{9n\ell^2}{2^{2h+1}}\right)
  >0.
\end{align*}
The last inequality is exactly the second hypothesis in
\eqref{eq:hypotheses}. Hence $c_\ell>0$, and $G$ contains a simple
cycle of length $\ell$.
Since $\ell$ was an arbitrary even integer satisfying
\eqref{eq:hypotheses}, this proves the asserted conclusion for every
such length. Finally, for $\ell>0$,
\[
  2^{2h}>\frac92\,n\ell^2
  \quad\Longleftrightarrow\quad
  \ell<\frac{\sqrt2}{3}\,\frac{2^h}{\sqrt n},
\]
which gives the equivalent interval formulation \eqref{eq:interval}.
\end{proof}

\bigskip
\begin{corollary}[The original girth condition]\label{cor:power-two}
Every finite simple cubic graph on $n\ge64$ vertices with girth
\[
  g\ge\log_2 n+2\log_2\!\bigl(\log_2 n\bigr)+9
\]
contains a simple cycle of every even length $\ell$ satisfying
\begin{equation}\label{eq:cor-interval}
  2\log_2 n+4\le\ell<\frac{16}{3}\log_2 n.
\end{equation}
In particular, it contains a cycle whose length is a power of two.
\end{corollary}

\begin{proof}
Put $t:=\log_2 n\ge6$ and
$h:=\lfloor(g-1)/2\rfloor$. The girth assumption and $2h\ge g-2$
give
\begin{equation}\label{eq:cor-girth}
  2^{2h}\ge2^{g-2}
  \ge2^{\,t+2\log_2 t+7}
  =128nt^2.
\end{equation}
For every even integer $\ell$ with $2t+4\le\ell<(16/3)t$,
\[
  \frac92\,n\ell^2
  <\frac92\,n\left(\frac{16}{3}t\right)^2
  =128nt^2
  \le2^{2h}.
\]
Both conditions of Theorem~\ref{thm:main} therefore hold, proving
\eqref{eq:cor-interval}.

To obtain a power-of-two length, let
\[
  q:=2^{\left\lceil\log_2(2t+4)\right\rceil},
\]
the least power of two not smaller than $2t+4$. Since
$u\le\lceil u\rceil<u+1$ for every real $u$,
\[
  2t+4\le q<4t+8\le\frac{16}{3}t,
\]
where the final inequality is equivalent to $t\ge6$.
Also $q\ge2t+4\ge16$, so $q$ is even.
Thus $q$ belongs to the interval in \eqref{eq:cor-interval}, and the
first part of the corollary supplies a simple cycle of length $q$.
\end{proof}

\end{document}
